\par\Needspace{6\baselineskip}\medskip
\hypertarget{native-218-article159}{}\pdfbookmark[2]{159.}{article-159}
\begin{center}159.\end{center}
Ad exemplum \emph{tertium} observationes sequentes Cereris eligimus, quarum
prima Bremae a clar. \textsc{Olbers}, secunda Gottingae a clar. \textsc{Harding}, tertia Lilienthalii a clar. \textsc{Bessel} instituta est:
\begin{center}\begingroup\footnotesize\setlength{\tabcolsep}{3pt}\renewcommand{\arraystretch}{1.3}
\begin{tabular}{|c|c|c|}
\hline
\shortstack{Tempus medium loci\\observationis} & Asc. recta & Declin. boreal. \\
\hline
1805 Sept. 5. $13^{\mathrm h}$ $8^{\mathrm m}$ $54^{\mathrm s}$ & $95^0\,59'\,25''$ & $22^0\,21'\,25''$ \\
1806 Ian. 17. $10$ $58$ $51$ & $101\quad18\quad40{,}6$ & $30\quad21\quad22{,}3$ \\
1806 Maii 23. $10$ $23$ $53$ & $121\quad56\quad7$ & $28\quad2\quad45$ \\
\hline
\end{tabular}
\endgroup\end{center}

Quum methodi, per quas parallaxis et aberrationis rationem habere licet,
si distantiae a terra tamquam omnino incognitae spectantur, per duo exempla
praecedentia abunde iam illustratae sint: superfluae laboris augmentationi in
hoc tertio exemplo renunciabimus, distantiasque approximatas e Commercio
litterario clar. de \textsc{Zach} (Vol. XI p. 284) eum in finem excerpemus, ut observationes ab effectu parallaxis et aberrationis purgentur. Has distantias una
cum reductionibus inde derivatis tabula sequens exhibet:
\begin{center}\begingroup\footnotesize\setlength{\tabcolsep}{3pt}\renewcommand{\arraystretch}{1.3}
\begin{tabular}{l|c|c|c}
Distantia Cereris a terra & $2{,}899$ & $1{,}638$ & $2{,}964$ \\
\shortstack[l]{tempus, intra quod lumen\\ad terram descendit} & $23^{\mathrm m}49^{\mathrm s}$ & $13^{\mathrm m}28^{\mathrm s}$ & $24^{\mathrm m}21^{\mathrm s}$ \\
tempus observationis reductum & $12^{\mathrm h}45^{\mathrm m}5^{\mathrm s}$ & $10^{\mathrm h}45^{\mathrm m}23^{\mathrm s}$ & $9^{\mathrm h}59^{\mathrm m}32^{\mathrm s}$ \\
tempus sidereum in gradibus & $355^0\,55'$ & $97^0\,59'$ & $210^0\,41'$ \\
parallaxis ascensionis rectae & $+1''{,}90$ & $+0''{,}22$ & $-1''{,}97$ \\
parallaxis declinationis & $-2{,}08$ & $-1{,}90$ & $-2{,}04$ \\
\end{tabular}
\endgroup\end{center}

\SourceLeaf{219}Problematis itaque data, postquam a parallaxi et aberratione liberata,
temporaque ad meridianum Parisinum reducta sunt, ita se habent:
\begin{center}\begingroup\footnotesize\setlength{\tabcolsep}{3pt}\renewcommand{\arraystretch}{1.3}
\begin{tabular}{|c|c|c|}
\hline
 & Asc. recta & Declinatio \\
\hline
1805 Sept. 5. $12^{\mathrm h}$ $19^{\mathrm m}$ $14^{\mathrm s}$ & $95^0\,59'\,23''{,}10$ & $22^0\,21'\,27''{,}08$ \\
1806 Ian. 17. $10$ $15$ $2$ & $101\quad18\quad40{,}38$ & $30\quad21\quad24{,}20$ \\
1806 Maii 23. $9$ $33$ $18$ & $121\quad56\quad8{,}97$ & $28\quad2\quad47{,}04$ \\
\hline
\end{tabular}
\endgroup\end{center}

Ex his ascensionibus rectis et declinationibus deductae sunt longitudines
et latitudines adhibita obliquitate eclipticae $23^0\,27'\,55''{,}90$, $23^0\,27'\,54''{,}59$,
$23^0\,27'\,53''{,}27$; dein longitudines a nutatione purgatae sunt, quae resp. fuit
$+17''{,}31$, $+17''{,}88$, $+18''{,}00$, posteaque ad initium anni 1806 reductae, applicata praecessione $+15''{,}98$, $-2''{,}39$, $-19''{,}68$. Denique pro temporibus reductis
e tabulis excerpta sunt loca Solis, ubi in longitudinibus nutatio praetermissa,
contra praecessio perinde ut longitudinibus Cereris adiecta est. Latitudo Solis
omnino neglecta. Hoc modo numeri sequentes in calculo adhibendi resultaverunt:
\begin{center}\begingroup\footnotesize\setlength{\tabcolsep}{3pt}\renewcommand{\arraystretch}{1.3}
\begin{tabular}{l|c|c|c}
Tempus 1805 Sept. & $5{,}51336$ & $139{,}42711$ & $265{,}39813$ \\
$a$, $a'$, $a''$ & $95^0\,32'\,18''{,}56$ & $99^0\,49'\,5''{,}87$ & $118^0\,5'\,28''{,}85$ \\
$\beta$, $\beta'$, $\beta''$ & $-0\quad59\quad34{,}06$ & $+7\quad16\quad36{,}80$ & $+7\quad38\quad49{,}39$ \\
$l$, $l'$, $l''$ & $342\quad54\quad56{,}00$ & $117\quad12\quad43{,}25$ & $241\quad58\quad50{,}71$ \\
$\log R$, $\log R'$, $\log R''$ & $0{,}003\,1514$ & $9{,}992\,9861$ & $0{,}005\,6974$ \\
\end{tabular}
\endgroup\end{center}

Iam calculi praeliminares in artt. 136---143 explicati sequentia suppeditant:
\begin{center}\begingroup\footnotesize\setlength{\tabcolsep}{3pt}\renewcommand{\arraystretch}{1.3}
\begin{tabular}{l|c|c|c}
$\gamma$, $\gamma'$, $\gamma''$ & $358^0\,55'\,28''{,}09$ & $156^0\,52'\,11''{,}49$ & $170^0\,48'\,44''{,}79$ \\
$\delta$, $\delta'$, $\delta''$ & $112\quad37\quad9{,}66$ & $18\quad48\quad39{,}81$ & $123\quad32\quad52{,}13$ \\
$A'D$, $AD'$, $AD''$ & $15\quad32\quad41{,}40$ & $252\quad42\quad19{,}14$ & $136\quad2\quad22{,}38$ \\
$A''D$, $A''D'$, $A'D''$ & $138\quad45\quad4{,}60$ & $6\quad26\quad41{,}10$ & $358\quad5\quad57{,}00$ \\
$\epsilon$, $\epsilon'$, $\epsilon''$ & $29\quad18\quad8{,}21$ & $170\quad32\quad59{,}08$ & $156\quad6\quad25{,}25$ \\
\end{tabular}
\endgroup\end{center}
\SourceLeaf{220}
\begin{center}\begingroup\footnotesize\setlength{\tabcolsep}{3pt}\renewcommand{\arraystretch}{1.3}
\begin{tabular}{ll}
$\sigma=8^0\,52'\,4''{,}05$ & \\
$\log a=0{,}184\,0193_n$, & $a=-1{,}527\,6340$\\
$\log b=0{,}004\,0987$ & \\
$\log\varkappa=0{,}161\,1012$ & \\
$\log\lambda=9{,}916\,4090_n$ &
\end{tabular}
\par\smallskip
\begin{tabular}{r@{\;=\;}l}
$\log c$ & $2{,}006\,6735$\\
$d$ & $117{,}50873$\\
$\log e$ & $0{,}856\,8244$\\
$\log\varkappa''$ & $9{,}977\,0819_n$\\
$\log\lambda''$ & $9{,}732\,0127_n$.
\end{tabular}
\endgroup\end{center}

Intervallum temporis inter observationem primam et secundam est dierum
$133{,}91375$, inter secundam et tertiam $125{,}97102$: hinc fit $\log\theta = 0{,}335\,8520$,
$\log\theta'' = 0{,}362\,4066$, $\log\dfrac{\theta''}{\theta} = 0{,}026\,5546$, $\log\theta\theta'' = 0{,}698\,2586$. Iam praecipua
momenta hypothesium trium primarum deinceps formatarum in conspectu sequenti exhibemus:
\begin{center}\begingroup\footnotesize\setlength{\tabcolsep}{3pt}\renewcommand{\arraystretch}{1.3}
\begin{tabular}{l|c|c|c}
 & I & II & III \\
\hline
$\log P = x$ & $0{,}026\,5546$ & $0{,}025\,6968$ & $0{,}025\,6275$ \\
$\log Q = y$ & $0{,}698\,2586$ & $0{,}739\,0190$ & $0{,}748\,1055$ \\
$\omega + \sigma$ & $7^0\,15'\,13''{,}523$ & $7^0\,14'\,47''{,}199$ & $7^0\,14'\,45''{,}071$ \\
$\log Qc\sin\omega$ & $1{,}154\,6650_n$ & $1{,}197\,3880_n$ & $1{,}206\,6327_n$ \\
$z$ & $7^0\,3'\,59''{,}018$ & $7^0\,2'\,32''{,}938$ & $7^0\,2'\,16''{,}900$ \\
$\log r'$ & $0{,}411\,4726$ & $0{,}412\,9371$ & $0{,}413\,2107$ \\
$\zeta$ & $160^0\,10'\,46''{,}74$ & $160^0\,20'\,7''{,}82$ & $160^0\,22'\,9''{,}42$ \\
$\zeta''$ & $262\quad6\quad1{,}03$ & $262\quad12\quad18{,}26$ & $262\quad14\quad19{,}49$ \\
$\log r$ & $0{,}432\,3934$ & $0{,}429\,1773$ & $0{,}428\,4841$ \\
$\log r''$ & $0{,}409\,4712$ & $0{,}407\,1975$ & $0{,}406\,4697$ \\
$\tfrac12(u''+u)$ & $262^0\,55'\,23''{,}22$ & $262^0\,57'\,6''{,}83$ & $262^0\,57'\,31''{,}17$ \\
$\tfrac12(u''-u)$ & $273\quad28\quad50{,}95$ & $273\quad29\quad15{,}06$ & $273\quad29\quad19{,}56$ \\
$2f'$ & $62\quad34\quad28{,}40$ & $62\quad49\quad56{,}50$ & $62\quad53\quad57{,}06$ \\
$2f$ & $31\quad8\quad30{,}03$ & $31\quad15\quad59{,}09$ & $31\quad18\quad13{,}83$ \\
$2f''$ & $31\quad25\quad58{,}45$ & $31\quad33\quad57{,}32$ & $31\quad35\quad43{,}32$ \\
$\log \eta$ & $0{,}020\,2496$ & $0{,}020\,3158$ & $0{,}020\,3494$ \\
$\log \eta''$ & $0{,}021\,1074$ & $0{,}021\,2429$ & $0{,}021\,2751$ \\
$\log P'$ & $0{,}025\,6968$ & $0{,}025\,6275$ & $0{,}025\,6289$ \\
$\log Q'$ & $0{,}739\,0190$ & $0{,}748\,1055$ & $0{,}750\,2337$ \\
$X$ & $-0{,}000\,8578$ & $-0{,}000\,0693$ & $+0{,}000\,0014$ \\
$Y$ & $+0{,}040\,7604$ & $+0{,}009\,0865$ & $+0{,}002\,1282$ \\
\end{tabular}
\endgroup\end{center}
\SourceLeaf{221}

Iam designando tres valores ipsius $X$ per $A$, $A'$, $A''$; tres valores ipsius $Y$
per $B$, $B'$, $B''$; quotientes e divisione quantitatum $A'B''-A''B'$, $A''B-AB''$,
$AB'-A'B$ per earundem aggregatum ortos resp. per $k$, $k'$, $k''$, ita ut habeatur $k+k'+k''=1$; denique valores ipsorum $\log P'$ et $\log Q'$ in hypothesi
tertia per $M$ et $N$ (qui forent valores novi ipsarum $x$, $y$, si hypothesin
quartam perinde e tertia derivare conveniret, ut tertia e secunda derivata
fuerat): e formulis art. 120 facile colligitur, valorem correctum ipsius $x$ fieri
$=M-k(A'+A'')-k'A''$, valoremque correctum ipsius $y = N-k(B'+B'')-k'B''$.
Calculo facto prior eruitur $=0{,}025\,6331$, posterior $=0{,}750\,9143$. Hisce valoribus correctis iam \emph{hypothesin quartam} superstruimus, cuius praecipua momenta haec sunt:
\begin{center}\begingroup\footnotesize\setlength{\tabcolsep}{3pt}\renewcommand{\arraystretch}{1.3}
\begin{tabular}{r|r}
$\omega + \sigma=7^0\,14'\,45''{,}247$ & $\log r''=0{,}406\,2033$\\
$\log Qc\sin\omega=1{,}209\,4284_n$ & $\tfrac12(u''+u)=262^0\,57'\,38''{,}78$\\
$z=7^0\,2'\,12''{,}736$ & $\tfrac12(u''-u)=273\quad29\quad20{,}73$\\
$\log r'=0{,}413\,2817$ & $2f'=62\quad55\quad16{,}64$\\
$\zeta=160^0\,22'\,45''{,}38$ & $2f=31\quad19\quad1{,}49$\\
$\zeta''=262\quad15\quad3{,}90$ & $2f''=31\quad36\quad15{,}20$\\
$\log r=0{,}428\,2792$ & \\
\end{tabular}
\endgroup\end{center}

Inter $2f'$ et $2f+2f''$ differentia $0''{,}05$ emergit, quam ita distribuemus, ut
statuamus $2f=31^0\,19'\,1''{,}47$, $2f''=31^0\,36'\,15''{,}17$. Quodsi iam e duobus locis
extremis elementa ipsa determinantur, sequentes numeri resultant:
\begin{center}\begingroup\footnotesize\setlength{\tabcolsep}{3pt}\renewcommand{\arraystretch}{1.3}
\begin{tabular}{ll}
\shortstack[l]{Anomalia vera pro loco primo\\.\,.\,.\,.\,.\,.\,.\,.\,.\,.\,.\,.\,.\,.\,.\,.} & $289^0\,7'\,39''{,}75$\\
\shortstack[l]{Anomalia vera pro loco tertio\\.\,.\,.\,.\,.\,.\,.\,.\,.\,.\,.\,.\,.\,.\,.\,.} & $352\quad2\quad56{,}39$\\
\shortstack[l]{Anomalia media pro loco primo\\.\,.\,.\,.\,.\,.\,.\,.\,.\,.\,.\,.\,.\,.\,.} & $297\quad41\quad35{,}65$\\
\shortstack[l]{Anomalia media pro loco tertio\\.\,.\,.\,.\,.\,.\,.\,.\,.\,.\,.\,.\,.\,.\,.} & $353\quad15\quad22{,}49$\\
\shortstack[l]{Motus medius diurnus sidereus\\.\,.\,.\,.\,.\,.\,.\,.\,.\,.\,.\,.\,.\,.\,.\,.\,.} & $769''{,}6750$\\
\shortstack[l]{Anomalia media pro initio anni 1806\\.\,.\,.\,.\,.\,.\,.\,.\,.\,.} & $322^0\,35'\,52''{,}51$\\
\shortstack[l]{Angulus $\varphi$\\.\,.\,.\,.\,.\,.\,.\,.\,.\,.\,.\,.\,.\,.\,.\,.\,.\,.\,.\,.\,.\,.\,.\,.\,.\,.\,.\,.\,.\,.\,.\,.\,.} & $4\quad37\quad57{,}78$\\
\shortstack[l]{Logarithmus semiaxis maioris\\.\,.\,.\,.\,.\,.\,.\,.\,.\,.\,.\,.\,.\,.\,.\,.\,.} & $0{,}442\,4661$.
\end{tabular}
\endgroup\end{center}

Computando ex hisce elementis locum heliocentricum pro tempore observationis mediae, invenitur anomalia media $326^0\,19'\,25''{,}72$, logarithmus radii
vectoris $0{,}413\,2825$, anomalia vera $320^0\,43'\,54''{,}87$: haecce distare deberet ab \SourceLeaf{222} anomalia vera pro loco primo differentia $2f''$, sive ab anomalia vera pro loco
tertio differentia $2f$, adeoque fieri deberet $=320^{0}\,43'\,54''{,}92$, sicuti
logarithmus radii vectoris $=0{,}413\,2817$: differentia $0''{,}05$ in anomalia vera,
octoque unitatum in isto logarithmo nullius momenti censenda est.

Si hypothesis quarta eodem modo ad finem perduceretur, ut tres praecedentes,
prodiret $X=0$, $Y=-0{,}000\,0168$, unde valores correcti ipsarum $x$, $y$ hi
colligerentur
\[
\begin{aligned}x &= \log P = 0{,}025\,6331 \\ &\quad\text{(idem ut in hypothesi quarta)}
\\
y &= \log Q = 0{,}750\,8973.
\end{aligned}
\]
\noindent Quibus valoribus si hypothesis quinta superstrueretur, solutio ultimam quam
tabulae permittunt praecisionem nancisceretur: sed elementa hinc resultantia vix
sensibiliter ab iis discreparent, quae hypothesis quarta suggessit.

Ut elementa completa habeantur, nihil iam superest, nisi ut situs plani orbitae
computetur. Ad normam praeceptorum art. 149 hic prodit
\begin{center}\begingroup\footnotesize\setlength{\tabcolsep}{3pt}\renewcommand{\arraystretch}{1.3}
\begin{tabular}{ll}\multicolumn{2}{c}{e loco primo}\\
\shortstack[l]{$g$\\.\,.\,.\,.\,.\,.\,.\,.\,.\,.\,.\,.\,.\,.\,.\,.\,.\,.\,.\,.\,.\,.\,.\,.\,.\,.\,.\,.} & $354^0\,9'\,44''{,}22$\\
\shortstack[l]{$h$\\.\,.\,.\,.\,.\,.\,.\,.\,.\,.\,.\,.\,.\,.\,.\,.\,.\,.\,.\,.\,.\,.\,.\,.\,.\,.\,.\,.} & $261\quad56\quad6{,}94$\\
\shortstack[l]{$i$\\.\,.\,.\,.\,.\,.\,.\,.\,.\,.\,.\,.\,.\,.\,.\,.\,.\,.\,.\,.\,.\,.\,.\,.\,.\,.\,.\,.} & $10\quad37\quad33{,}02$\\
\shortstack[l]{$\AscendingNode$\\.\,.\,.\,.\,.\,.\,.\,.\,.\,.\,.\,.\,.\,.\,.\,.\,.\,.\,.\,.\,.\,.\,.\,.\,.\,.\,.} & $80\quad58\quad49{,}06$\\
Distantia perihelii a nodo ascendente & $65\quad2\quad4{,}47$\\
\shortstack[l]{Longitudo perihelii\\.\,.\,.\,.\,.\,.\,.\,.\,.\,.\,.\,.\,.} & $146\quad0\quad53{,}53$\\
\end{tabular}
\par\smallskip
\begin{tabular}{ll}\multicolumn{2}{c}{e loco tertio}\\
\shortstack[l]{$g''$\\.\,.\,.\,.\,.} & $57^0\,5'\,0''{,}91$\\
\shortstack[l]{$h''$\\.\,.\,.\,.\,.} & $161\quad0\quad1{,}61$\\
.\,.\,.\,.\,.\,.\,. & $10\quad37\quad33{,}00$\\
.\,.\,.\,.\,.\,.\,. & $80\quad58\quad49{,}10$\\
.\,.\,.\,.\,.\,.\,. & $65\quad2\quad4{,}52$\\
.\,.\,.\,.\,.\,.\,. & $146\quad0\quad53{,}62$.
\end{tabular}
\endgroup\end{center}
\noindent Sumto itaque medio statuetur $i=10^{0}\,37'\,33''{,}01$, $\AscendingNode=80^{0}\,58'\,49''{,}08$,
longitudo perihelii $=146^{0}\,0'\,53''{,}57$. Denique longitudo media pro initio
anni 1806 erit $=108^{0}\,36'\,46''{,}08$.
